% Appendix: LabVIEW Quick Reference
% Auto-generated from megn300_master_reference.tex

\chapter*{Appendix: LabVIEW Quick Reference for MEGN 300}
\addcontentsline{toc}{chapter}{Appendix: LabVIEW Quick Reference for MEGN 300}

\begin{figure}[htbp]
\centering
\begin{tikzpicture}[
    step/.style={draw=MinesBlue, fill=LightBG, thick, rounded corners=3pt,
                 minimum width=3.0cm, minimum height=1.0cm, align=center, font=\scriptsize\sffamily\bfseries},
    loop/.style={draw=AccentTeal, fill=AccentTeal!8, thick, rounded corners=3pt,
                 minimum width=3.0cm, minimum height=1.0cm, align=center, font=\scriptsize\sffamily\bfseries},
    arr/.style={-{Stealth[length=2.5mm]}, thick, MinesBlue},
    note/.style={font=\tiny\sffamily\itshape, text=MinesSilver, align=left, anchor=west},
]
\node[step] (create) at (0,0) {DAQmx\\Create Channel};
\node[step] (timing) at (0,-1.8) {DAQmx Timing};
\node[step] (start) at (0,-3.6) {DAQmx Start Task};
\node[loop] (read) at (0,-5.4) {DAQmx Read\\(inside While Loop)};
\node[step] (stop) at (0,-7.2) {DAQmx Stop Task};
\node[step] (clear) at (0,-9.0) {DAQmx Clear Task};

\draw[arr] (create) -- (timing);
\draw[arr] (timing) -- (start);
\draw[arr] (start) -- (read);
\draw[arr] (read) -- (stop);
\draw[arr] (stop) -- (clear);

% Loop indicator
\draw[AccentTeal, thick, dashed, rounded corners=5pt] (-2,-4.8) rectangle (2,-6.0);
\node[font=\tiny\sffamily\bfseries, AccentTeal] at (1.5,-4.6) {While Loop};

% Notes
\node[note] at (2.2,0) {Channel type, physical channel,\\terminal config (Diff or RSE)};
\node[note] at (2.2,-1.8) {Sample rate, samples/channel,\\continuous vs. finite};
\node[note] at (2.2,-3.6) {Arms the acquisition\\(call once, outside loop)};
\node[note] at (2.2,-5.4) {Read data each iteration\\Process, display, log here};
\node[note] at (2.2,-7.2) {Stop acquisition\\(after loop exits)};
\node[note] at (2.2,-9.0) {Release hardware resource\\(\textbf{omitting this = ``resource reserved''})};

% Error wire
\draw[WarnOrange, thick, decorate, decoration={snake, amplitude=1pt, segment length=5pt}]
    (1.5,-0.4) -- (1.5,-8.6);
\node[font=\tiny\sffamily\bfseries, WarnOrange, rotate=90, anchor=south] at (1.8,-4.5) {Error wire (thread all VIs)};
\end{tikzpicture}
\caption{The standard DAQmx programming pattern for MEGN~300. Every LabVIEW VI that reads from the NI DAQ should follow this sequence. The error wire (orange) must connect through every VI to ensure proper cleanup on failure.}
\label{fig:daqmx_pattern}
\end{figure}

This appendix provides a condensed reference for the LabVIEW programming patterns most commonly used in MEGN~300 labs.

\section*{Core Concepts}
LabVIEW programs are called \textbf{Virtual Instruments (VIs)}. Each VI has a \textbf{Front Panel} (the user interface: controls for input, indicators for output) and a \textbf{Block Diagram} (the code: wired graphical nodes that execute dataflow). Data flows from left to right along wires; a node executes when all its inputs are available.

\section*{DAQ Configuration}
\begin{enumerate}[leftmargin=1.5em]
\item Open NI~MAX (Measurement \& Automation Explorer). Verify your DAQ device appears under ``Devices and Interfaces.''
\item Right-click the device $\to$ ``Test Panels'' to verify analog input/output and digital I/O without writing any code.
\item In LabVIEW, use the \textbf{DAQ Assistant Express VI} for quick setup, or the \textbf{DAQmx VIs} for full control (recommended for MEGN~300 labs).
\end{enumerate}

\textbf{DAQmx programming pattern}:
\begin{enumerate}[nosep,leftmargin=1.5em]
\item \texttt{DAQmx Create Virtual Channel} (specify physical channel, measurement type, range).
\item \texttt{DAQmx Timing} (set sample rate and samples per channel).
\item \texttt{DAQmx Start Task}.
\item \textbf{Loop}: \texttt{DAQmx Read} $\to$ process data $\to$ display/log.
\item \texttt{DAQmx Stop Task} $\to$ \texttt{DAQmx Clear Task}.
\end{enumerate}

\section*{Common Programming Patterns}

\subsection*{While Loop with Stop Button}
The fundamental LabVIEW pattern: a While Loop on the Block Diagram with a Boolean ``Stop'' button on the Front Panel wired to the loop's conditional terminal. All DAQ reads, processing, and display happen inside the loop. The loop runs until the user presses Stop.

\subsection*{Shift Registers for Running Calculations}
Shift registers (right-click the loop border $\to$ ``Add Shift Register'') pass data from one iteration to the next. Use them for: running averages, PID integral accumulation, previous-value storage (for derivative calculations), and event counting.

\subsection*{Timing with Wait (ms)}
Place a ``Wait (ms)'' function inside the While Loop to control the loop rate. For a 100\,ms loop period (10\,Hz): wire a constant ``100'' to the Wait input. Without this, the loop runs as fast as possible, consuming CPU and potentially overwhelming the DAQ.

\subsection*{File I/O for Data Logging}
Use the \textbf{Write to Measurement File Express VI} for simple logging, or \textbf{Write to Spreadsheet File} for custom formats. Always include a timestamp column. Log to a file \emph{inside} the While Loop for real-time data capture; do not accumulate data in memory and write at the end (you lose everything if LabVIEW crashes).

\section*{Signal Processing VIs}

\begin{table}[htbp]\centering\small
\begin{tabularx}{\textwidth}{l X l}
\toprule
\textbf{VI / Function} & \textbf{Purpose} & \textbf{Location} \\
\midrule
FFT Power Spectrum & Compute single-sided amplitude or power spectrum & Signal Processing $\to$ Spectral \\
Filter (Express VI) & Apply LP/HP/BP/notch filter with interactive design & Signal Processing $\to$ Filters \\
IIR/FIR Filter Design & Programmatic filter design with full control & Signal Processing $\to$ Filters \\
Mean / Std Dev & Basic statistics on a waveform & Mathematics $\to$ Probability \\
Curve Fit & Polynomial or custom equation fit to data & Mathematics $\to$ Fitting \\
PID (autotuning) & PID controller with autotuning option & Control Design $\to$ PID \\
Waveform Chart & Real-time scrolling display & Graph Indicators \\
XY Graph & Plot any X vs.\ Y data (e.g., calibration curves) & Graph Indicators \\
Build Array / Index & Construct and access array elements & Array Functions \\
Bundle / Unbundle & Group multiple data types into a cluster & Cluster Functions \\
\bottomrule
\end{tabularx}
\caption{Essential LabVIEW VIs for MEGN~300.}
\end{table}

\section*{PID Controller Implementation}
LabVIEW provides a built-in PID VI (\texttt{PID.vi}) in the Control Design toolkit. Inputs: setpoint, process variable, PID gains ($K_p$, $K_i$, $K_d$), output range, and dt. Outputs: control output and PID components (P, I, D individually for debugging).

\textbf{Implementation pattern}:
\begin{enumerate}[nosep,leftmargin=1.5em]
\item Read the process variable from the DAQ (temperature, speed, etc.).
\item Wire the setpoint (Front Panel control) and PV to PID.vi inputs.
\item Wire the PID output to the actuator (DAQ analog output for motor speed, or digital output for SSR duty cycle).
\item Place PID.vi inside the While Loop with a consistent loop rate (e.g., 100\,ms).
\item Add Front Panel controls for $K_p$, $K_i$, $K_d$ so you can tune in real time.
\item Add a Waveform Chart showing setpoint, PV, and control output vs.\ time.
\end{enumerate}

\begin{tipbox}[Debugging LabVIEW VIs]
\begin{itemize}[nosep,leftmargin=1.5em]
\item \textbf{Highlight Execution} (light bulb icon): animates dataflow so you can see values moving along wires. Slow, but reveals logic errors.
\item \textbf{Probes}: right-click any wire $\to$ ``Probe'' to display its value in real time without adding indicators.
\item \textbf{Breakpoints}: pause execution at any node to inspect data.
\item \textbf{Broken wire = broken VI}: LabVIEW will not run until all wires are valid. Follow the broken wire to the source error.
\item \textbf{Error clusters}: wire the ``error in'' and ``error out'' terminals of DAQmx VIs in series. If any step fails, downstream VIs skip execution and propagate the error message to a ``Simple Error Handler'' at the end.
\end{itemize}
\end{tipbox}


\section*{My VI is Broken: The Signal Walk Debugging Protocol}

When your measurement or control system produces incorrect results, do \textbf{not} change code randomly. Walk the signal chain from physical input to software output, verifying each stage:

\begin{enumerate}[leftmargin=1.5em]
\item \textbf{Stage 1: Verify the raw sensor voltage with a multimeter.} Disconnect the DAQ. Measure the sensor output directly. If the voltage is wrong here, the problem is hardware (wiring, sensor, excitation), not software.
\item \textbf{Stage 2: Verify the DAQ is reading correctly.} Reconnect the DAQ. In NI~MAX Test Panels, read the analog input channel. Compare to the multimeter reading. If they disagree, check: channel number, input mode (Differential vs.\ RSE), input range, and wiring polarity.
\item \textbf{Stage 3: Verify the LabVIEW data path.} Add Probes (right-click wire $\to$ Probe) at key points: after the DAQ Read, after the calibration equation, after the filter, and at the PID input. Compare each value to what you expect.
\item \textbf{Stage 4: Verify the actuator responds.} Force the PID output to a manual 50\% (bypass the PID, wire a constant). Does the heater/motor respond? If not, the problem is the output circuit (MOSFET, SSR, wiring), not the control algorithm.
\item \textbf{Stage 5: Close the loop.} Only after stages 1--4 pass should you troubleshoot PID tuning. If the sensor reads correctly, the actuator responds correctly, and the loop still oscillates, the gains need adjustment (reduce $K_p$ first).
\end{enumerate}

This systematic isolation prevents the most common debugging failure: ``poking at the code'' to fix a hardware wiring error.

\section*{Common Mistakes in MEGN 300 Labs}
\begin{enumerate}[leftmargin=1.5em]
\item \textbf{No loop timing}: the While Loop runs at maximum speed, overwhelming the DAQ buffer and producing erratic sampling rates. Always add a Wait (ms) function.
\item \textbf{DAQ task not closed}: forgetting \texttt{DAQmx Clear Task} leaves the DAQ in a locked state. The next run fails with ``resource reserved.'' Always wire Clear Task outside the loop (executes on exit).
\item \textbf{Wrong channel configuration}: selecting ``RSE'' (referenced single-ended) when the sensor requires ``Differential'' mode. This adds noise and ground offset errors.
\item \textbf{Waveform Chart vs.\ Waveform Graph}: the Chart plots data point-by-point as it arrives (real-time display). The Graph plots an entire array at once (post-acquisition). Use Chart inside the loop; Graph outside.
\item \textbf{Incorrect PID loop rate}: the PID VI uses $dt$ in its integral and derivative calculations. If the loop rate varies (no timing control), the PID gains effectively change from iteration to iteration. Use a consistent, controlled loop period.
\item \textbf{Forgetting anti-aliasing}: sampling at 1\,kS/s without a hardware anti-alias filter means any noise above 500\,Hz aliases into your measurement. The NI~USB-6001 has no built-in anti-alias filter; you must add an external RC or active filter.
\end{enumerate}

% ================================================================
